\section{PLANTEAMIENTO DEL PROBLEMA}

    Según \parencite{bib:Serway2014}, la termodinámica es parte de la física clásica (desarrollada desde antes de 1900), estudia las relaciones entre el calor y su transformación a otras formas de energía.

    La \textbf{Psicrometría} es la rama de la termodinámica que estudia las propiedades de las mezclas gaseosas (gas y vapor); las mezcla del aire seco y vapor de agua (conocido com aire húmedo) es la la mezcla más estudiada debido a su importancia en la vida diaria, su uso y aplicación se extiende desde la industria química hasta la meteorología, ya que permite entender cómo se comporta el aire en diferentes condiciones de temperatura y humedad.

    para entender la psicrometría, es necesario conocer algunos conceptos básicos de la termodinámica, como la temperatura, la presión de la mezcla de compuestos en estado gaseoso.



    \subsection*{CONCEPTOS}
    Son diversas los conceptos que se han manejado en el estudio de la psicrometría; a continuación se mencionan algunos de ellos.

    Para \parencite[ p. 25]{bib:SamuelC.Sugarman2007} la psicrometría es la ciencia de la termodinámica del aire húmedo aplicado a los sistemas \glsxtrshort{hvac}. 
    % Enfriamento: Disminución de la temperature de aire.
    % Teoría de procesos de humidifcación y deshumidificación


    % \textcite{Rajput2007}, hace una distición entre los términos, psicrometría y psicrométrica; donde la psicrometría conceptualiza como el arte   

    % Los diagramas que se han revisado 

    % De acuerdo a  \href{www.wikipedia.org}{wikipedia} 
    Para \textcite[]{bib:Himmelblau2004}, un diagrama de humedad o carta psicrométrica, no es mas que una representación gráfica del balance de materia y energía de la mezcla vapor de agua y aire seco; evidentemente asociados a las variables psicrométricas.


    Según \textcite[]{bib:Martinez1992}, la psicrometría es el estudio del aire húmedo. Mientras que \textcite[]{bib:cengel2015termodinamica} evita usar el término psicrometría y solamente refiere como mezclas de gas-vapor y acondicionamiento de aire.    




    % Para () señala que diversos autores definen la psicrometría como la ciencia que estudia las propiedades de la mezcla de aire (se entiende aire por aire seco) y vapor de agua; en el ámbito de procesos ()

    Para la elaboración de este proyecto, se tomará la definición de \parencite[]{bib:Perry2000}, donde la psicrometría se refiere a la determinación de las propiedades de la mezcla de gas-vapor, donde la mezcla aire seco-vapor de agua es el sistema más representativo.



    % La psicrometría es una rama de la ciencias físicas que se ocupa de estudiar las propiedades termodinámicas de mezclas de gas-vapor; en particular, el aire húmedo; sin embargo, puede extenderse a mezclas de otros compuestos.

    % Dentro del area de ingeniería (HVAC) su uso está bastante extendido, el cual se ocupa en dar soluciones de un entorno confortable el cual consiste en:


